% Subsection 2.1.1: Transformer 与自回归生成机制

Transformer 架构是当前主流代码大语言模型的基础网络结构\cite{vaswani2017attention,devlin2019bert}，其核心通过多头自注意力机制完成序列依赖建模，具备优秀的长距离特征捕捉与并行计算能力。在代码生成场景中，模型通常采用解码器-only（纯解码器）结构，以自回归方式逐 token 生成程序语句，能够较好地适配代码的语法结构、逻辑关系与执行流程。

自回归生成的基本原理是：在给定已生成前缀序列的条件下，依次预测下一个位置的最优 token，直至生成结束标记。该过程可表示为条件概率最大化问题：
\[
P(x_t \mid x_1, x_2, \dots, x_{t-1}; \theta)
\]
其中 \(x_t\) 为当前待预测的 token，\(x_1,\dots,x_{t-1}\) 为已生成序列，\(\theta\) 表示模型参数。在解码阶段，可通过贪心搜索、束搜索或随机采样等策略确定输出序列，分别在输出确定性、多样性与正确性之间取得不同权衡。

得益于自注意力机制对全局上下文的建模能力，Transformer 能够同时关注问题描述、约束条件与历史生成内容，适合算法代码所需的逻辑一致性与结构规范性。图~\ref{fig:ch2-autoregressive} 给出该自回归生成过程的逐步示意：每一步均以「问题描述 + 已生成前缀」为输入，经 Transformer 解码器输出下一个 token 的条件概率分布，采样后追加回前缀，直至生成 \texttt{[EOS]}。但自回归范式存在误差累积效应：早期生成的错误会沿序列不断扩散，最终导致逻辑错误、边界处理缺失或结构异常，这也是算法代码生成稳定性不足的重要原因。

\begin{figure}[!htb]
  \centering
  \begin{tikzpicture}[
    node distance=4mm and 5mm,
    box/.style={draw, rounded corners=2pt, align=center, minimum height=7mm, minimum width=14mm, font=\footnotesize},
    model/.style={draw, fill=black!5, rounded corners=2pt, align=center, minimum height=10mm, minimum width=18mm, font=\footnotesize},
    tok/.style={draw, rounded corners=1pt, align=center, minimum height=6mm, minimum width=10mm, font=\scriptsize\ttfamily},
    arr/.style={-{Stealth[length=2mm]}, thick},
  ]
    \node[box] (q) {问题描述\\ \scriptsize prompt $x_{1{:}p}$};
    \node[tok, right=of q] (t1) {def};
    \node[tok, right=of t1] (t2) {solve};
    \node[tok, right=of t2] (t3) {(n)};
    \node[font=\footnotesize, right=2mm of t3] (dots) {$\cdots$};
    \node[tok, right=2mm of dots] (tk) {[EOS]};

    \node[model, below=8mm of t2] (m) {Transformer\\Decoder $\theta$};
    \draw[arr] (q.south) to [bend right=15] (m.west);
    \draw[arr] (t1.south) -- (m.north -| t1);
    \draw[arr] (t2.south) -- (m.north);
    \draw[arr] (t3.south) -- (m.north -| t3);

    \node[font=\scriptsize, below=2mm of m] (p) {$P(x_t \mid x_{1:t-1};\theta)$};
    \draw[arr] (m.east) to [bend right=10] node[midway, right, font=\scriptsize]{采样 $x_t$} (tk.south);
    \draw[arr] (tk.west) to [bend right=20] node[midway, above, font=\scriptsize]{追加至前缀} (t3.east);
  \end{tikzpicture}
  \caption{Transformer 解码器的自回归生成：每步以「问题 + 已生成前缀」为条件预测下一 token，循环直至 \texttt{[EOS]}}
  \label{fig:ch2-autoregressive}
\end{figure}